\begin{proof}[Proof of \cref{obj:prop:unrestricted-near-complete-phase-boundary}]
\begin{enumerate}
\item For all sufficiently large \(n\), we have
\[
n\ge1,
\qquad
1-q_n\le\frac{M}{\sqrt n}.
\]
Fix such an \(n\). Since \(d_n\ge1\) and \(0<q_n\le1\), \cref{obj:thm:unrestricted-near-complete-arrival-envelope} applies. Retaining the last term in its upper-bound minimum gives
\[
\frac{1}{256n}
\le \mathfrak R^{+}_{n,d_n,q_n}
\le \frac4n+(1-q_n)^2.
\]
% lean: CausalSmith.Stat.MarRareqLogfrontier.unrestricted_near_complete_arrival_envelope

\item Since \(n>0\), we have \(\sqrt n>0\), so multiplying the deficit bound by \(\sqrt n\) yields
\[
(1-q_n)\sqrt n\le M.
\]
Moreover, \(q_n\le1\) implies
\[
0\le(1-q_n)\sqrt n.
\]
Thus, using \(M\ge0\), we may square the inequality to obtain
\[
\bigl((1-q_n)\sqrt n\bigr)^2\le M^2.
\]
The identity \((\sqrt n)^2=n\) then gives
\[
n(1-q_n)^2
=
\bigl((1-q_n)\sqrt n\bigr)^2
\le M^2.
\]
Dividing by \(n>0\), we conclude that
\[
(1-q_n)^2\le\frac{M^2}{n}.
\]
% lean: CausalSmith.Stat.MarRareqLogfrontier.unrestricted_near_complete_phase_boundary

\item Substituting this bound into the upper estimate gives
\[
\mathfrak R^{+}_{n,d_n,q_n}
\le \frac4n+(1-q_n)^2
\le \frac4n+\frac{M^2}{n}
=\frac{4+M^2}{n}.
\]
Together with the lower estimate, this proves both asserted inequalities for all sufficiently large \(n\).
% lean: CausalSmith.Stat.MarRareqLogfrontier.unrestricted_near_complete_phase_boundary
\end{enumerate}
\end{proof}
